\section{Test Cases}

\begin{table}[htbp]
    \centering
    \caption{Planned test-case ladder for the manuscript.}
    \begin{tabular}{llll}
        \toprule
        Stage & Geometry source & Role & Planned output \\
        \midrule
        1D & tracked 1D benchmark & analytical limit & line/profile PNG \\
        2D & 4StraightFractures & transparent image case & 2D panel PNG \\
        3D & Kelvin1-200x200x200 & simple 3D topology & slice/render PNG \\
        3D large & Est\_3phase500cubed4micron & porous-media demo & multi-panel PNG \\
        \bottomrule
    \end{tabular}
\end{table}

\subsection{Simple 2D Image-Based Case}

The planned 2D case reuses the existing segmented image asset:
\begin{itemize}
\item \texttt{images/raw/4StraightFractures.raw.tar.gz}
\item metadata: \texttt{images/raw/4StraightFractures.txt}
\end{itemize}

This case should play the same role as the simple 2D case in the reference
paper: a geometrically transparent test that makes the coupling between
transport, internal surface interaction, and retardation easy to interpret.

Planned figure set:
\begin{itemize}
\item geometry PNG,
\item aqueous concentration PNG,
\item surface loading PNG.
\end{itemize}

The 2D case should be written to answer three questions:
\begin{enumerate}
\item does the method recover the expected directional retardation pattern in a
simple geometry,
\item does the surface-loading field track the geometry in an interpretable
way,
\item does the 2D image case remain visually and numerically easier to audit
than the 3D cases.
\end{enumerate}

\subsection{Simple 3D Image-Based Case}

The planned simple 3D case reuses:
\begin{itemize}
\item \texttt{images/raw/Kelvin1-200x200x200.raw.tar.gz}
\item metadata: \texttt{images/raw/Kelvin1-200x200x200.txt}
\end{itemize}

This case should demonstrate the transition from 2D transport pathways to a
fully 3D unresolved surface network while still remaining easy to interpret.

The simple 3D section should focus on:
\begin{itemize}
\item how 3D topology changes the transport path,
\item how the unresolved internal surface distribution changes the uptake
pattern,
\item how much the 3D case increases chemistry cost relative to the 2D case.
\end{itemize}

\subsection{Large 3D Porous-Media Case}

The planned large 3D porous-media case reuses:
\begin{itemize}
\item \texttt{images/raw/Est\_3phase500cubed4micron.raw.tar.gz}
\item metadata: \texttt{images/raw/Est\_3phase500cubed4micron.txt}
\end{itemize}

This case is the application-scale demonstration. It should use an
Estaillades-style multi-scale image so the method is exercised on a domain with
more realistic multi-scale structural contrast than the controlled 2D and
simple 3D cases.

The large porous-media section should report:
\begin{itemize}
\item mesh/image size,
\item effective resolution,
\item visible multi-scale structural contrast in the reported slices,
\item MPI rank count,
\item representative runtime,
\item one concentration metric and one surface-loading metric that remain
comparable to the earlier sections.
\end{itemize}

\begin{figure}[htbp]
    \centering
    \fbox{\rule{0pt}{2in}\rule{0.9\linewidth}{0pt}}
    \caption{Placeholder for the 2D case figure set
    (planned PNG target: \texttt{figures/png/fig02\_2d\_fractures.png}).}
\end{figure}

\begin{figure}[htbp]
    \centering
    \fbox{\rule{0pt}{2in}\rule{0.9\linewidth}{0pt}}
    \caption{Placeholder for the simple 3D case figure set
    (planned PNG target: \texttt{figures/png/fig03\_3d\_kelvin.png}).}
\end{figure}

\begin{figure}[htbp]
    \centering
    \fbox{\rule{0pt}{2in}\rule{0.9\linewidth}{0pt}}
    \caption{Placeholder for the large 3D porous-media Estaillades case figure
    set (planned PNG target: \texttt{figures/png/fig04\_3d\_estaillades.png}).}
\end{figure}
